\subsection{Detection of Small Lesions on the MSLesSeg Test Set}
\label{sec:main_results}

\Cref{tab:loss_comparison_small_lesion} summarizes the results on the six
held-out MSLesSeg test patients (12 scans), and \cref{tab:paired_differences}
gives the paired differences between CATMIL and each baseline with
patient-level 95\% CIs (\cref{sec:stats}).

\begin{table}[t]
\centering
\caption{Comparison of CATMIL with standard losses and the instance-level
loss CC-DiceCE on the MSLesSeg test set, without post-processing. Values
are the mean $\pm$ sample standard deviation over the five fold models.
Best values are shown in \textbf{bold}; second-best values are
\underline{underlined} (ranked before rounding).}
\label{tab:loss_comparison_small_lesion}
\resizebox{\textwidth}{!}{%
\begin{tabular}{lccccc}
\toprule
\textbf{Metric}
& \textbf{CATMIL (ours)}
& \textbf{DiceCE}
& \textbf{Tversky}
& \textbf{FocalTversky}
& \textbf{CC-DiceCE} \\
\midrule
Small Lesion Recall $\uparrow$
& \textbf{0.873 $\pm$ 0.01}
& 0.796 $\pm$ 0.03
& \underline{0.834 $\pm$ 0.02}
& 0.831 $\pm$ 0.02
& 0.831 $\pm$ 0.01 \\
FN Count $\downarrow$
& \textbf{2.57 $\pm$ 0.15}
& 4.95 $\pm$ 0.40
& 3.77 $\pm$ 0.23
& 3.70 $\pm$ 0.24
& \underline{3.65 $\pm$ 0.33} \\
FN Volume Fraction $\downarrow$
& \textbf{0.021 $\pm$ 0.003}
& 0.034 $\pm$ 0.01
& 0.026 $\pm$ 0.01
& \underline{0.025 $\pm$ 0.005}
& 0.035 $\pm$ 0.01 \\
\midrule
DSC $\uparrow$
& \textbf{0.783 $\pm$ 0.01}
& 0.780 $\pm$ 0.01
& 0.771 $\pm$ 0.01
& \underline{0.780 $\pm$ 0.01}
& 0.778 $\pm$ 0.01 \\
HD95 (mm) $\downarrow$
& \textbf{7.98 $\pm$ 1.84}
& 9.04 $\pm$ 1.34
& 10.24 $\pm$ 0.15
& \underline{8.21 $\pm$ 1.93}
& 9.72 $\pm$ 1.53 \\
Lesion F1 $\uparrow$
& 0.757 $\pm$ 0.02
& \underline{0.843 $\pm$ 0.01}
& 0.840 $\pm$ 0.01
& \textbf{0.846 $\pm$ 0.01}
& 0.806 $\pm$ 0.01 \\
FP Volume (mm$^3$) $\downarrow$
& \underline{1537 $\pm$ 245}
& 1819 $\pm$ 654
& 2621 $\pm$ 559
& 2282 $\pm$ 710
& \textbf{1514 $\pm$ 289} \\
\bottomrule
\end{tabular}%
}
\end{table}

\textbf{CATMIL detects more small lesions than every baseline.} Small-lesion
recall was 0.873 with CATMIL, against 0.796 with DiceCE and 0.831--0.834
with the other losses. The gain was present in every test patient against
every baseline, and its 95\% CI excluded zero in each comparison
(\cref{tab:paired_differences}). Weighting each lesion within an overlap
term, as CC-DiceCE does, therefore does not by itself recover the lesions
that CATMIL recovers. The gain over DiceCE and CC-DiceCE also held when at least half of a
lesion's voxels had to be covered to count as detected, whereas the margin
over the Tversky-based losses shrank to about 0.01 under this criterion
(\cref{appendix:strict-detection}).

\textbf{CATMIL halves the number of completely missed lesions.} CATMIL
missed 2.6 lesions per scan, against 5.0 with DiceCE (about 48\% fewer) and
3.7--3.8 with the other losses. The reduction held in all six patients
against every baseline except FocalTversky (five of six).

\begin{table}[t]
\centering
\caption{Paired differences between CATMIL and each baseline (CATMIL minus
baseline) for the main detection metrics, with 95\% patient-level bootstrap
CIs and, in parentheses, the number of test patients in which CATMIL is
better. Differences whose CI excludes zero are shown in \textbf{bold}. PP5:
after removing predicted components smaller than 5 voxels. All metrics are
given in \cref{tab:paired_differences_full}.}
\label{tab:paired_differences}
\resizebox{\textwidth}{!}{%
\begin{tabular}{lllll}
\toprule
\textbf{Metric} & \textbf{vs.\ DiceCE} & \textbf{vs.\ Tversky} & \textbf{vs.\ FocalTversky} & \textbf{vs.\ CC-DiceCE} \\
\midrule
\multicolumn{5}{l}{\emph{MSLesSeg}} \\
Small Lesion Recall $\uparrow$ & \textbf{+0.077 [+0.030, +0.157] (6/6)} & \textbf{+0.039 [+0.004, +0.093] (6/6)} & \textbf{+0.042 [+0.005, +0.101] (6/6)} & \textbf{+0.042 [+0.020, +0.071] (6/6)} \\
FN Count $\downarrow$ & \textbf{$-$2.38 [$-$5.29, $-$0.84] (6/6)} & \textbf{$-$1.20 [$-$2.98, $-$0.27] (6/6)} & \textbf{$-$1.13 [$-$2.71, $-$0.27] (5/6)} & \textbf{$-$1.08 [$-$1.95, $-$0.52] (6/6)} \\
Lesion F1 $\uparrow$ & \textbf{$-$0.086 [$-$0.138, $-$0.043] (0/6)} & \textbf{$-$0.083 [$-$0.143, $-$0.038] (0/6)} & \textbf{$-$0.088 [$-$0.151, $-$0.039] (0/6)} & \textbf{$-$0.049 [$-$0.091, $-$0.018] (0/6)} \\
Lesion F1, PP5 $\uparrow$ & $-$0.017 [$-$0.040, +0.001] (2/6) & $-$0.012 [$-$0.038, +0.006] (2/6) & $-$0.017 [$-$0.043, +0.001] (2/6) & +0.003 [$-$0.010, +0.017] (4/6) \\
\midrule
\multicolumn{5}{l}{\emph{3D-MR-MS}} \\
Small Lesion Recall $\uparrow$ & \textbf{+0.081 [+0.034, +0.127] (5/6)} & \textbf{+0.054 [+0.019, +0.086] (5/6)} & \textbf{+0.055 [+0.021, +0.087] (5/6)} & +0.039 [$-$0.001, +0.080] (4/6) \\
FN Count $\downarrow$ & \textbf{$-$5.13 [$-$12.40, $-$0.50] (5/6)} & \textbf{$-$4.43 [$-$11.23, $-$0.27] (5/6)} & \textbf{$-$4.43 [$-$11.47, $-$0.23] (5/6)} & \textbf{$-$2.27 [$-$4.80, $-$0.47] (5/6)} \\
Lesion F1 $\uparrow$ & \textbf{$-$0.197 [$-$0.258, $-$0.142] (0/6)} & \textbf{$-$0.168 [$-$0.242, $-$0.106] (0/6)} & \textbf{$-$0.166 [$-$0.232, $-$0.108] (0/6)} & \textbf{$-$0.131 [$-$0.195, $-$0.066] (0/6)} \\
\bottomrule
\end{tabular}%
}
\end{table}

\textbf{Overlap and false-positive volume are not worsened.} The DSC and
HD95 of CATMIL were comparable to those of the baselines, and its
false-positive volume was similar to that of DiceCE and CC-DiceCE and lower
than that of the Tversky-based losses (\cref{tab:paired_differences_full}).
The additional detections are therefore not obtained by predicting more
lesion tissue.

\textbf{The cost is a lower lesion-wise precision.} The lesion-wise F1 of
CATMIL was 0.757, against 0.806--0.846 for the baselines, and lower in every
patient. Together with the low false-positive volume, this points to
numerous small false-positive components rather than over-segmentation
(\cref{sec:fp_pp}).

\subsection{Detection of Lesions of at Least 3~mm}
\label{sec:clinical_size}

The small-lesion threshold of 150 voxels used above is a technical
definition. \Cref{tab:clinical_size} therefore reports recall for lesions
above and below the clinical size threshold of 3~mm in longest diameter
(\cref{sec:intro}).

\begin{table}[t]
\centering
\caption{Lesion recall by lesion size relative to the 3~mm longest
diameter used in clinical MRI reading. $n$: number of reference lesions.
Recall is the fraction of lesions detected, averaged over the five fold
models. $\Delta$: CATMIL minus DiceCE or CC-DiceCE with 95\% patient-level
bootstrap CI; differences whose CI excludes zero are shown in
\textbf{bold}.}
\label{tab:clinical_size}
\resizebox{\textwidth}{!}{%
\begin{tabular}{llrcccccll}
\toprule
\textbf{Dataset} & \textbf{Lesions} & $n$ & \textbf{CATMIL} & \textbf{DiceCE} & \textbf{Tversky} & \textbf{FocalTversky} & \textbf{CC-DiceCE} & $\Delta$ vs.\ DiceCE & $\Delta$ vs.\ CC-DiceCE \\
\midrule
MSLesSeg & $\geq$3~mm & 345 & 0.944 & 0.870 & 0.907 & 0.907 & 0.910 & \textbf{+0.074 [+0.032, +0.129]} & \textbf{+0.035 [+0.017, +0.057]} \\
         & $\geq$3~mm and $\leq$150 voxels & 233 & 0.924 & 0.827 & 0.874 & 0.875 & 0.884 & \textbf{+0.098 [+0.040, +0.171]} & \textbf{+0.040 [+0.020, +0.066]} \\
         & $<$3~mm & 16 & 0.275 & 0.088 & 0.175 & 0.225 & 0.213 & \textbf{+0.188 [+0.044, +0.300]} & +0.063 [$-$0.100, +0.188] \\
\midrule
3D-MR-MS & $\geq$3~mm & 269 & 0.894 & 0.830 & 0.837 & 0.833 & 0.863 & \textbf{+0.065 [+0.015, +0.125]} & \textbf{+0.031 [+0.010, +0.096]} \\
         & $\geq$3~mm and $\leq$150 voxels & 156 & 0.855 & 0.753 & 0.768 & 0.756 & 0.817 & \textbf{+0.103 [+0.050, +0.203]} & \textbf{+0.038 [+0.010, +0.129]} \\
         & $<$3~mm & 120 & 0.393 & 0.282 & 0.300 & 0.308 & 0.350 & \textbf{+0.112 [+0.010, +0.167]} & +0.043 [$-$0.013, +0.075] \\
\bottomrule
\end{tabular}%
}
\end{table}

\textbf{The detection gain holds for lesions of clinical size.} For lesions
of at least 3~mm, CATMIL detected 94.4\% of the reference lesions on
MSLesSeg, against 87.0\% for DiceCE and 91.0\% for CC-DiceCE, and 89.4\% on
3D-MR-MS, against 83.0\% and 86.3\%; all four differences excluded zero.
The largest gain was for small lesions just above the clinical size (at
least 3~mm and at most 150 voxels), for which CATMIL raised recall over
DiceCE by about 10 percentage points on both datasets. These are lesions
that a reader would count but that a voxel-wise objective tends to miss.
Lesions below 3~mm were rare on MSLesSeg (16 of 361) but common on 3D-MR-MS
(120 of 389), and their recall remained low for every loss (at most 0.39);
some of them may still be relevant in specific
locations~\cite{barkhof_2024_2025}, but lesion location was not analyzed.

\subsection{Which Lesions Are Recovered}

\begin{figure}[t]
    \centering
    \includegraphics[width=0.85\textwidth]{figures/recall_vs_size_curve.pdf}
    \caption{Lesion recall against lesion size on the MSLesSeg test set.
    Upper panel: fraction of reference lesions detected in each size
    interval (logarithmic axis); lines are the mean over the five fold
    models and shaded bands span the fold minimum to maximum. Lower panel:
    number of reference lesions per interval. Dotted vertical lines mark
    10, 50, and 200~mm$^3$. A lesion counts as detected if any predicted
    voxel overlaps it.}
    \label{fig:recall_vs_size_curve}
\end{figure}

\textbf{CATMIL moves the detection limit toward smaller lesions.} For
lesions of 5--30~mm$^3$, CATMIL detected 87\% of lesions, against 58\% for
DiceCE and 70--80\% for the other losses (\cref{fig:recall_vs_size_curve}).
It reached a recall of 0.8 at about 10~mm$^3$, against about 18--40~mm$^3$ for
the baselines. Above about 100~mm$^3$ all losses detect nearly every lesion,
and lesions of at most 5~mm$^3$ remain difficult for every loss.

\textbf{The gain is a net gain.} Pairing
every reference lesion between CATMIL and each baseline trained on the same
fold (\cref{fig:lesion_fate}), CATMIL detected 30.6 lesions per fold model
that DiceCE missed, while DiceCE detected only 2.0 that CATMIL missed (15.4
against 2.4 for CC-DiceCE), and no lesion above 200~mm$^3$ was lost.
CATMIL missed fewer lesions than DiceCE and CC-DiceCE in 10 of the 12 test
scans and the same number in the other two; against Tversky and
FocalTversky, it missed more lesions in one and two scans, respectively
(\cref{fig:fn_per_case_paired}).
The gain was concentrated in scans with many small lesions
(\cref{tab:test_case_lesion_profile}): in P47\_T1, CATMIL recovered about 11
lesions relative to DiceCE while the scan-level DSC remained 0.85, because
the recovered lesions are too small to change a volume-overlap score.

\begin{figure}[t]
    \centering
    \includegraphics[width=\textwidth]{figures/lesion_fate.pdf}
    \caption{Lesions detected by only one of two losses on the MSLesSeg
    test set, by lesion size. Each panel pairs every reference lesion
    between CATMIL and one baseline trained on the same fold. Bars to the
    right: lesions detected by CATMIL but not by the baseline; bars to the
    left: lesions detected by the baseline but not by CATMIL. Grey text:
    lesions detected by both or by neither. Values are means over the five
    fold models.}
    \label{fig:lesion_fate}
\end{figure}

\begin{figure}[t]
    \centering
    \includegraphics[width=\textwidth]{figures/fn_per_case_paired.pdf}
    \caption{Missed lesions per test case for CATMIL compared with DiceCE,
    Tversky, FocalTversky, and CC-DiceCE. Each point is one test case, placed at the
    mean FN lesion count over the five fold models; error bars span the
    fold minimum to maximum. Points below the dashed diagonal (shaded
    region) are cases in which CATMIL misses fewer lesions than the
    baseline. The text in each panel gives the number of cases in which
    CATMIL misses fewer, the same number, or more lesions, and the mean
    number of missed lesions per case for the baseline and for CATMIL.}
    \label{fig:fn_per_case_paired}
\end{figure}

\textbf{The baselines do not respond to most of the lesions they miss.}
Among lesions of at most 10~mm$^3$, the maximum predicted probability stayed
below 0.1 for 82\% under DiceCE, 62--68\% under the other baselines, and
47\% under CATMIL; intermediate responses (0.1--0.5) were rare for every
loss (\cref{fig:max_prob_by_size}). These lesions therefore cannot be
recovered by lowering the decision threshold of an existing model: CATMIL
makes the network respond to lesions that a voxel-wise objective leaves
invisible.

\begin{figure}[t]
    \centering
    \includegraphics[width=0.75\textwidth]{figures/max_prob_by_size.pdf}
    \caption{Maximum foreground probability inside each reference lesion on
    the MSLesSeg test set, by lesion size. Bars give the fraction of
    lesions whose maximum probability is below 0.1 (no response), between
    0.1 and 0.5 (weak response below the threshold), or at least 0.5
    (detected). Each lesion contributes once per fold model; $n$ is the
    number of reference lesions in the bin.}
    \label{fig:max_prob_by_size}
\end{figure}

\subsection{Lesion Examples}

\Cref{fig:qual_lesion_gallery} shows individual lesions in 24~mm crops,
selected by a fixed rule applied to the detection results of all five fold
models rather than by visual inspection (\cref{appendix:gallery-rule}).

\begin{figure}[!htbp]
    \centering
    \includegraphics[width=\textwidth]{figures/qual_lesion_gallery.pdf}
    \caption{Example lesions from the MSLesSeg test set, selected by a fixed
    rule (\cref{appendix:gallery-rule}). Each row shows one lesion or one
    CATMIL false-positive component: the axial slice with the crop location
    (cyan box), a 24~mm FLAIR crop with the target (cyan arrow), the
    reference contour (white), and the fold-0 prediction of each loss.
    Green: correctly segmented lesion voxels; red: missed lesion voxels;
    yellow: false-positive voxels. Bottom left of each prediction panel:
    number of the five fold models that detect the lesion; bottom right
    ($p$): maximum predicted probability inside the lesion, averaged over
    the five models. Only FLAIR is shown, although the models also receive
    T1-WI and T2-WI.}
    \label{fig:qual_lesion_gallery}
\end{figure}

\textbf{Recovered lesions (rows 1--4).} For three of the four recovered
lesions, the maximum DiceCE probability inside the lesion was at most 0.01,
whereas CATMIL assigned 0.93--0.99 and detected each lesion in all five
models. The lesions illustrate different difficulties: a 5-voxel lesion
3~mm from a larger neighbor (row~1); a clearly visible but isolated 13-voxel
lesion (row~2), which shows that the omission is an effect of size rather
than of visibility; a faint 21-voxel lesion next to a large lesion that all
losses detect (row~3); and a 37-voxel juxtacortical lesion that has to be
separated from the bright cortex on FLAIR (row~4). Juxtacortical lesions
contribute to dissemination in space~\cite{montalban_diagnosis_2025}, so
recovering them matters beyond the lesion count. CATMIL located these
lesions but outlined them only partly (lesion-wise DSC 0.28--0.52), and in
row~1 it joined the lesion to its neighbor: it mainly corrects the omission
of lesions rather than their delineation.

\textbf{Remaining misses and false positives (rows 5--6).} The 18-voxel
posterior fossa lesion in row~5 was missed by all twenty models. Its
contrast is lower than that of 99\% of the small test lesions, and the
earlier scan of the same patient has no lesion within 5~mm, so it is either
new at follow-up or an annotation that the earlier scan does not support.
Across the test set, the lesions of 5--50 voxels that CATMIL still missed
had a median FLAIR contrast of 1.86, against 3.65 for those it detected:
the remaining misses are limited by image evidence rather than by the
training objective. The 3-voxel CATMIL false positive in row~6 is a small
focal hyperintensity; it may be a tiny unannotated lesion, a vessel, or
noise, and it is removed by the size filter described next.

\subsection{False Positives and a Minimum-Size Filter}
\label{sec:fp_pp}

At its default output, CATMIL predicted about 51 components per scan,
against about 30 reference lesions (lesion-level recall 0.909, precision
0.674), while its false-positive volume was close to the lowest among the
compared losses (1537~mm$^3$, against 1514~mm$^3$ for CC-DiceCE). The false positives are therefore numerous small components, most of
them close to lesions: the median distance of false-positive voxels to the
nearest reference lesion is 1--2~mm.

Removing predicted components smaller than 5 voxels (PP5;
\cref{tab:loss_comparison_small_lesion_pp5}) raised the lesion-wise F1 of
CATMIL from 0.757 to 0.841, close to that of the baselines, while its
small-lesion recall decreased only from 0.873 to 0.855 and remained the
highest. After PP5, the lesion-wise F1 difference was no longer supported
against any baseline (\cref{tab:paired_differences}), whereas the
small-lesion recall gain remained in all six patients against every
baseline (\cref{tab:paired_differences_full}). On the larger sample of the
cross-validation patients, a small lesion-wise F1 cost remained after PP5
(\cref{sec:catmil_oof}).

\begin{table}[t]
\centering
\caption{Comparison of CATMIL with standard losses and the instance-level
loss CC-DiceCE on the MSLesSeg test set after post-processing with a
minimum connected-component size of 5 voxels (PP5). Best values are shown
in \textbf{bold}; second-best values are \underline{underlined} (ranked
before rounding).}
\label{tab:loss_comparison_small_lesion_pp5}
\resizebox{\textwidth}{!}{%
\begin{tabular}{lccccc}
\toprule
\textbf{Metric}
& \textbf{CATMIL (ours)}
& \textbf{DiceCE}
& \textbf{Tversky}
& \textbf{FocalTversky}
& \textbf{CC-DiceCE} \\
\midrule
Small Lesion Recall $\uparrow$
& \textbf{0.855 $\pm$ 0.01}
& 0.785 $\pm$ 0.03
& \underline{0.822 $\pm$ 0.02}
& 0.820 $\pm$ 0.02
& 0.815 $\pm$ 0.01 \\
FN Count $\downarrow$
& \textbf{3.13 $\pm$ 0.21}
& 5.23 $\pm$ 0.52
& \underline{4.07 $\pm$ 0.24}
& 4.08 $\pm$ 0.24
& 4.18 $\pm$ 0.30 \\
FN Volume Fraction $\downarrow$
& \textbf{0.024 $\pm$ 0.004}
& 0.035 $\pm$ 0.01
& 0.028 $\pm$ 0.01
& \underline{0.027 $\pm$ 0.01}
& 0.037 $\pm$ 0.01 \\
\midrule
DSC $\uparrow$
& \textbf{0.784 $\pm$ 0.01}
& 0.780 $\pm$ 0.01
& 0.771 $\pm$ 0.01
& \underline{0.780 $\pm$ 0.01}
& 0.778 $\pm$ 0.01 \\
HD95 (mm) $\downarrow$
& \textbf{8.10 $\pm$ 1.76}
& 9.09 $\pm$ 1.37
& 10.36 $\pm$ 0.21
& \underline{8.23 $\pm$ 1.95}
& 9.82 $\pm$ 1.52 \\
Lesion F1 $\uparrow$
& 0.841 $\pm$ 0.01
& \textbf{0.859 $\pm$ 0.01}
& 0.853 $\pm$ 0.01
& \underline{0.858 $\pm$ 0.005}
& 0.839 $\pm$ 0.01 \\
FP Volume (mm$^3$) $\downarrow$
& \underline{1515 $\pm$ 244}
& 1815 $\pm$ 654
& 2618 $\pm$ 558
& 2279 $\pm$ 710
& \textbf{1505 $\pm$ 290} \\
\bottomrule
\end{tabular}%
}
\end{table}

When the threshold and filter size were tuned so that CATMIL matched the
lesion-wise precision of each baseline, it still detected more small
lesions than DiceCE and CC-DiceCE on MSLesSeg, but not than the
Tversky-based losses; with the same settings applied to 3D-MR-MS, the gain
remained against CC-DiceCE, whereas the CI against DiceCE narrowly included
zero (\cref{sec:matched_precision}).

\subsection{Evaluation on a Second Dataset: 3D-MR-MS}

\begin{table}[t]
\centering
\caption{Comparison of CATMIL with standard losses and the instance-level
loss CC-DiceCE on the 3D-MR-MS test set. Best values are shown in
\textbf{bold}; second-best values are \underline{underlined}.}
\label{tab:external_eval_3dmrms}
\resizebox{\textwidth}{!}{%
\begin{tabular}{lccccc}
\toprule
\textbf{Metric} & \textbf{CATMIL (ours)} & \textbf{DiceCE} & \textbf{Tversky} & \textbf{FocalTversky} & \textbf{CC-DiceCE} \\
\midrule
Small Lesion Recall $\uparrow$
& \textbf{0.618 $\pm$ 0.03}
& 0.537 $\pm$ 0.01
& 0.564 $\pm$ 0.02
& 0.563 $\pm$ 0.02
& \underline{0.579 $\pm$ 0.01} \\
FN Count $\downarrow$
& \textbf{16.9 $\pm$ 0.64}
& 22.0 $\pm$ 0.94
& 21.3 $\pm$ 1.15
& 21.3 $\pm$ 1.20
& \underline{19.1 $\pm$ 1.06} \\
FN Volume Fraction $\downarrow$
& 0.085 $\pm$ 0.03
& \textbf{0.068 $\pm$ 0.01}
& 0.082 $\pm$ 0.03
& \underline{0.074 $\pm$ 0.03}
& 0.113 $\pm$ 0.03 \\
\midrule
DSC $\uparrow$
& 0.718 $\pm$ 0.01
& \underline{0.730 $\pm$ 0.02}
& 0.726 $\pm$ 0.01
& \textbf{0.734 $\pm$ 0.01}
& 0.720 $\pm$ 0.02 \\
HD95 (mm) $\downarrow$
& \textbf{11.55 $\pm$ 1.27}
& \underline{11.94 $\pm$ 1.79}
& 13.38 $\pm$ 3.65
& 12.27 $\pm$ 3.10
& 13.91 $\pm$ 3.05 \\
Lesion F1 $\uparrow$
& 0.471 $\pm$ 0.04
& \textbf{0.668 $\pm$ 0.02}
& \underline{0.639 $\pm$ 0.02}
& 0.637 $\pm$ 0.03
& 0.602 $\pm$ 0.03 \\
FP Volume (mm$^3$) $\downarrow$
& \textbf{1232 $\pm$ 165}
& \underline{1234 $\pm$ 213}
& 1994 $\pm$ 284
& 1958 $\pm$ 269
& 1283 $\pm$ 189 \\
\bottomrule
\end{tabular}%
}
\end{table}

On 3D-MR-MS (\cref{tab:external_eval_3dmrms}), CATMIL again had the
highest small-lesion recall (0.618, against 0.537 for DiceCE and 0.579 for
CC-DiceCE) and missed the fewest lesions (16.9 per scan, against 22.0 for
DiceCE, about 23\% fewer, and 19.1 for CC-DiceCE). Against the voxel-level
losses, both gains were present in five of the six test patients with CIs
that excluded zero (\cref{tab:paired_differences}); against CC-DiceCE,
CATMIL missed fewer lesions, whereas the recall gain narrowly included
zero. The false-positive volume was similar to that of DiceCE, and CATMIL
again detected far more lesions that DiceCE missed than the reverse (about
33 against 2.4 per fold model; \cref{appendix:lesion-level-analyses}).

The cost of the higher sensitivity was larger on this dataset. The
lesion-wise F1 of CATMIL was 0.471, against 0.602--0.668, and lower in every
test patient: without filtering, CATMIL produced about 150 components per
scan that overlapped no reference lesion, against 14--19 for the
voxel-level losses and 43 for CC-DiceCE, so a filter larger than PP5 is
needed here (\cref{appendix:lesion-level-analyses}). The DSC was also
slightly lower (0.718 against 0.730 for DiceCE, with a CI that included
zero). One of the six test patients contributes 249 of the 389 reference
lesions, so pooled lesion counts on this dataset are dominated by this
patient, whereas the patient-level analysis gives each patient equal
weight.

\subsection{Evaluation on the Cross-Validation Patients}
\label{sec:catmil_oof}

Each test set contains six patients. To test whether the MSLesSeg results
extend beyond them, the cross-validation predictions of the 81 scans of the
other 47 MSLesSeg patients were evaluated with the same metrics. Each scan
was predicted once, by the fold model that was not trained on it; the
validation folds were used only for monitoring during training, and the
CATMIL weights were fixed before the experiments. Paired differences were
assessed with a patient-level bootstrap over the 47 patients (2{,}000
resamples, reduced from the 10{,}000 used elsewhere to limit computation
over the 81 scans; \cref{tab:oof_paired}).

\begin{table}[t]
\centering
\caption{Paired differences between CATMIL and each baseline on the
cross-validation predictions of the 81 MSLesSeg training scans
(47 patients). Values are CATMIL minus baseline with the 95\% patient-level
bootstrap CI in brackets; differences whose CI excludes zero are shown in
\textbf{bold}. Lesion F1 is reported without filtering and after removing
components smaller than 5 voxels (PP5).}
\label{tab:oof_paired}
\resizebox{\textwidth}{!}{%
\begin{tabular}{lccccc}
\toprule
\textbf{CATMIL minus} & \textbf{Small Lesion Recall} & \textbf{FN Count} & \textbf{Lesion F1} & \textbf{Lesion F1 (PP5)} & \textbf{DSC} \\
\midrule
DiceCE       & \textbf{+0.110 [+0.074, +0.140]} & \textbf{$-$3.37 [$-$4.58, $-$2.35]} & \textbf{$-$0.101 [$-$0.135, $-$0.067]} & \textbf{$-$0.025 [$-$0.044, $-$0.007]} & +0.004 [$-$0.002, +0.010] \\
Tversky      & \textbf{+0.070 [+0.037, +0.099]} & \textbf{$-$2.06 [$-$2.90, $-$1.34]} & \textbf{$-$0.100 [$-$0.137, $-$0.065]} & \textbf{$-$0.023 [$-$0.041, $-$0.005]} & $-$0.001 [$-$0.007, +0.005] \\
FocalTversky & \textbf{+0.073 [+0.039, +0.100]} & \textbf{$-$2.23 [$-$3.13, $-$1.45]} & \textbf{$-$0.097 [$-$0.128, $-$0.064]} & \textbf{$-$0.026 [$-$0.043, $-$0.009]} & $-$0.003 [$-$0.008, +0.003] \\
CC-DiceCE    & \textbf{+0.054 [+0.026, +0.077]} & \textbf{$-$1.65 [$-$2.32, $-$1.11]} & \textbf{$-$0.078 [$-$0.106, $-$0.050]} & \textbf{$-$0.026 [$-$0.043, $-$0.009]} & $-$0.002 [$-$0.009, +0.004] \\
\bottomrule
\end{tabular}%
}
\end{table}

\textbf{The detection gain is reproduced on the larger sample.} CATMIL
raised small-lesion recall from 0.692 with DiceCE to 0.802 and reduced
missed lesions from 7.51 to 4.14 per scan. The gains against Tversky,
FocalTversky, and CC-DiceCE also excluded zero, with narrower intervals
than on the test set, while the DSC did not differ from that of any
baseline.

\textbf{The precision cost is also reproduced.} Without filtering, the
lesion-wise F1 of CATMIL was 0.078 to 0.101 lower than that of the
baselines; after PP5, the difference shrank to 0.023--0.026 but remained
below zero against every baseline, including CC-DiceCE. The equal
lesion-wise F1 of CATMIL and CC-DiceCE after PP5 on the test set is
therefore not reproduced.

\textbf{The objective matters more than the architecture.} Four further
architectures trained with the unchanged nnU-Net objective (SegResNet,
UMambaBot, SwinUNETR, and UNETR) did not detect more small lesions than
nnU-Net with the same objective on the test set, and CATMIL exceeded each of them in
small-lesion recall on both the test set and the cross-validation patients
(\cref{sec:arch_vs_objective}).
